\documentclass[tikz,border=10pt]{standalone}
\usepackage{fontspec}
\usepackage{xeCJK}
\setCJKmainfont[
  Path=/Users/YW/Library/Fonts/,
  Extension=.ttf
]{Microsoft Yahei}
\usepackage{tikz}
\usetikzlibrary{arrows.meta,positioning,calc}

\begin{document}
\begin{tikzpicture}[
  >=Stealth,
  signal/.style={->, line width=1.2pt},
  block/.style={draw, rectangle, minimum width=2.3cm, minimum height=1.0cm, line width=1.2pt},
  sum/.style={draw, circle, minimum size=0.55cm, inner sep=0pt, line width=1.2pt},
  tap/.style={draw, circle, minimum size=0.18cm, inner sep=0pt, fill=white, line width=1.2pt},
  note/.style={font=\small, align=left}
]

\node[font=\bfseries] at (4.8, 4.2) {方框图基本组件试重建};

\draw[signal] (0,3.1) -- (2.6,3.1);
\node[note, anchor=west] at (3.0,3.1) {信号线\\有箭头的直线，箭头表示信号传递方向};

\draw[signal] (0,1.5) -- (2.6,1.5);
\node[tap] (tap1) at (1.3,1.5) {};
\draw[signal] (tap1) -- ++(0,-0.9);
\node[note, anchor=west] at (3.0,1.5) {引出点\\从同一信号线上引出的信号，数值和性质完全相同};

\node[sum] (sum1) at (1.3,-0.2) {$\times$};
\draw[signal] (0.1,-0.2) -- (sum1);
\draw[signal] (1.3,0.9) -- (sum1);
\draw[signal] (sum1) -- (2.6,-0.2);
\node at (0.9,0.18) {$+$};
\node at (1.68,0.18) {$-$};
\node[note, anchor=west] at (3.0,-0.2) {综合点\\对两个或两个以上的信号进行代数运算};

\node[block] (g1) at (1.3,-2.3) {$G(s)$};
\draw[signal] (-0.1,-2.3) -- (g1.west);
\draw[signal] (g1.east) -- (2.7,-2.3);
\node[note, anchor=west] at (3.0,-2.3) {方框\\表示典型环节或其组合，框内写传递函数};

\end{tikzpicture}
\end{document}
